\documentclass[11pt]{article}

% ----------------------------------------------------------------------
% Packages
% ----------------------------------------------------------------------
\usepackage[utf8]{inputenc}
\usepackage[T1]{fontenc}
\usepackage[a4paper,margin=1in]{geometry}
\usepackage{booktabs}
\usepackage{xcolor}
\usepackage{listings}
\usepackage{enumitem}
\usepackage{microtype}
\usepackage{titlesec}
\usepackage[hidelinks]{hyperref}
\hypersetup{
  colorlinks=true,
  linkcolor=black,
  urlcolor=blue!60!black,
  citecolor=black
}

% ----------------------------------------------------------------------
% Listing style for JSON
% ----------------------------------------------------------------------
\definecolor{codebg}{gray}{0.97}
\definecolor{codekey}{rgb}{0.10,0.30,0.60}
\lstdefinestyle{json}{
  backgroundcolor=\color{codebg},
  basicstyle=\ttfamily\scriptsize,
  breaklines=true,
  columns=fullflexible,
  frame=single,
  rulecolor=\color{gray!40},
  showstringspaces=false,
  keepspaces=true,
  literate=
    *{0}{{{\color{codekey}0}}}{1}
     {:}{{{\color{gray}:}}}{1}
}

\newcommand{\code}[1]{\texttt{\small #1}}

\titleformat{\section}{\normalfont\large\bfseries}{\thesection}{0.6em}{}
\setlist{topsep=2pt,itemsep=1pt}

% ----------------------------------------------------------------------
% Title
% ----------------------------------------------------------------------
\title{\bfseries Pragmatics-Preserving Machine Translation\\ for Indian Dialogue\\[4pt]
\large Progress Report}
\author{Lohith Kola} % edit as needed
\date{\today}

\begin{document}
\maketitle

\begin{abstract}
\noindent
I am building a benchmark to measure whether machine translation systems preserve
the \emph{social} meaning of Indian conversational dialogue, and not only its
literal content. A translation can be word-for-word accurate and still be wrong:
if a student politely asks a professor for more time and the system renders it as a
blunt demand, no fact has changed, but the student now sounds rude. I call this a
pragmatic failure, and current evaluation practice does not measure it. This report
covers what I have done so far this summer. I began by fixing the research
direction, reviewing seven papers, and writing a design blueprint, and over the past
few weeks I turned that blueprint into a working instrument: a machine-readable
annotation schema, a draft pilot dataset of 47 dialogue items in Hindi and Hinglish,
a complete set of annotation and adjudication guidelines, and a validation and
analysis toolchain with an automated test suite. I want to be clear at the outset
that nothing in the dataset has yet been validated by a human, and I have been
careful not to manufacture any result that would require one. What I have produced
is the measuring apparatus. Putting it in front of native speakers is what comes
next.
\end{abstract}

% ======================================================================
\section{Introduction and motivation}
% ======================================================================

Machine translation is usually judged on whether it preserves what was said. The
question I am asking is different: does it preserve what was \emph{meant}, socially?
Indian conversation is a good place to ask this, because Hindi marks respect
grammatically and pervasively. A speaker cannot address someone without choosing
between the respectful \code{aap}, the familiar \code{tum}, and the intimate
\code{tu}, and that choice carries information that English has no direct way to
express. Every translation into English is therefore forced to make a social
decision that the source made explicitly, and any system can get that decision
wrong while still being perfectly accurate about the facts.

My working definition is deliberately narrow so that it can be measured:

\begin{quote}
\itshape
A pragmatic failure in machine translation occurs when a translation preserves the
basic semantic content of an utterance but fails to preserve its intended social or
communicative meaning in context.
\end{quote}

The primary contribution of this project is a benchmark and an evaluation protocol,
not a new translation model. My guiding research question is: to what extent do
current machine translation systems preserve the intended pragmatic and social
meaning of Indian conversational dialogue, and can that failure be measured
reliably enough to compare systems and interventions?

% ======================================================================
\section{What the literature told me}
% ======================================================================

I grounded the design in seven papers, which fall into three groups. The first
group shows that stylistic properties of translation are controllable: Sennrich et
al.~\cite{sennrich2016} steer a politeness distinction with a side constraint, and
Niu et al.~\cite{niu2017} and Niu and Carpuat~\cite{niu2020} control formality by
reranking and by synthetic supervision. The second group establishes that standard
quality scores miss context-dependent meaning. Bawden et al.~\cite{bawden2018}
build contrastive test suites and show that translation quality can rise while
discourse accuracy stays at chance. Voita et al.~\cite{voita2019a} find that around
7\% of English-Russian sentence pairs are individually good yet wrong when read
together, with forms of address the single largest source, and their repair work
\cite{voita2019b} shows the same problem can be fixed after the fact. The third
group is about evaluation itself: Agrawal et al.~\cite{agrawal2024} show that adding
context to a metric helps mainly for short and ambiguous turns, and that corrupted
context actively hurts, which is a caution against assuming more context is always
better.

Three lessons shaped everything that followed. First, semantic quality and social
appropriateness have to be scored separately, because a translation can be strong
on one and fail the other. Second, contrastive minimal pairs are the right tool:
if the wrong candidate is as semantically plausible as the right one, a system that
ignores social meaning is reduced to guessing. Third, human judgement stays
primary, and I should expect only moderate agreement on the hardest dimensions.
Niu and Carpuat report human agreement on formality at around 0.5 on
Krippendorff's alpha, on a task narrower than mine, so I am treating agreement in
the 0.4 to 0.6 range on the hard dimensions as expected rather than as failure.

I also identified the gap the project fills. None of the seven papers jointly cover
speech acts, indirect requests and refusals, the teasing-versus-insult boundary,
warmth, or discourse-motivated code-switching, and none work on Indian or
code-mixed data. That combination is where I am placing the benchmark.

% ======================================================================
\section{Benchmark design}
% ======================================================================

\subsection{Pragmatic categories}
I organised the benchmark around four categories, the first three primary and the
fourth included but treated as the most uncertain:

\begin{enumerate}[label=(\arabic*)]
  \item Politeness and formality
  \item Indirect requests and refusals
  \item Stance and emotion
  \item Discourse-motivated code-switching
\end{enumerate}

\subsection{The item}
Each item is a short dialogue of zero to three previous turns followed by the
utterance to be translated, annotated with speaker and listener roles, the social
relationship, a full set of pragmatic labels, one or more acceptable reference
translations, and a single contrastive translation that is semantically close but
pragmatically wrong. Listing~\ref{lst:item} shows an abridged item. The full schema
has 43 fields.

\begin{lstlisting}[style=json,caption={An abridged benchmark item (HIN\_POL\_001).},label={lst:item}]
{
  "item_id": "HIN_POL_001",
  "language": "Hindi",
  "context": [
    {"turn_id": 1, "speaker_role": "Professor",
     "text": "Yeh assignment kal shaam tak jama kar dijiye."}
  ],
  "source_utterance": {"turn_id": 2, "speaker_role": "Student",
     "text": "Sir, kya mujhe do din aur mil sakte hain?"},
  "literal_gloss": "Sir, can I get two more days?",
  "relationship": "ACADEMIC", "relative_status": "SPEAKER_LOWER",
  "primary_phenomenon": "POLITENESS_FORMALITY",
  "speech_act": "REQUEST", "politeness_level": "RESPECTFUL",
  "indirectness": "MODERATELY_INDIRECT",
  "context_status": "HELPFUL",
  "reference_translations": ["Sir, could I please have two more days?"],
  "contrastive_translation": "I need two more days.",
  "contrastive_error": "A deferential request becomes a flat demand to a superior.",
  "contrastive_error_category": "POLITENESS_ERROR", "severity": "MAJOR",
  "review_status": "NEEDS_NATIVE_REVIEW"
}
\end{lstlisting}

The reference translation ``Sir, could I please have two more days?'' and the
contrastive ``I need two more days'' say the same thing about the facts. What
separates them is that a respectful request has become a demand. Rating the first
as pragmatically correct and the second as incorrect, while rating both as
semantically adequate, is exactly the judgement the benchmark is built to elicit.

\subsection{Contrastive construction}
A contrastive translation is only useful if it fails for the right reason. I wrote
eight requirements for a valid one: it must preserve the propositional content, be
grammatical, be plausible read in isolation, become inappropriate only in context,
target a single pragmatic feature, avoid unrelated semantic errors, stay minimally
different from the reference, and be confirmed by native speakers. The most
dangerous failure mode is a contrastive that is actually still acceptable, because
then the item measures nothing and, worse, penalises a system for translating well.
I made catching that the single most valuable check I ask reviewers to perform.

\subsection{Context necessity and severity}
Two design choices are worth calling out. Each item records how much its context
matters, on a three-way scale of \code{NOT\_REQUIRED}, \code{HELPFUL}, and
\code{REQUIRED}, where \code{REQUIRED} means the intended reading genuinely changes
without the previous turns. This is what lets me later separate a real context
benefit from general system capability. Each contrastive error also carries a
severity of \code{MINOR}, \code{MAJOR}, or \code{CRITICAL}, which feeds a
severity-weighted error score so that a system making a few reversals of meaning is
distinguished from one making many cosmetic slips.

% ======================================================================
\section{Language scope}
% ======================================================================

I set the first release to Hindi and naturally occurring Hinglish, both translated
into English, and I keep Hindi and Hinglish as two distinct values rather than
merging them, so that I can report on monolingual and code-mixed input separately.
I chose this setting for practical and linguistic reasons: it gives the largest
available pool of native-speaker annotators, it is heavily used in exactly the
settings the benchmark models, and Hindi's obligatory politeness marking makes the
target phenomena frequent and concrete. All source text is written in romanized
Latin script, which is how Hindi is typically typed in chat and which lets Hindi
and Hinglish be compared directly, at the known cost that many systems are trained
on Devanagari.

I want to be explicit about one limitation, because it would be easy for the
project title to imply otherwise: this choice is not representative of Indian
languages in general. Hindi is one language among many with different politeness
systems, and I describe my findings as findings about Hindi and Hinglish, never as
findings about ``Indian language MT''. Telugu is planned as a second setting, but
only after the schema has been validated on one language pair, and the binding
constraint there is annotator availability rather than schema work.

% ======================================================================
\section{What I built}
% ======================================================================

Over the past few weeks I turned the design into a working repository.
Table~\ref{tab:built} lists the components.

\begin{table}[h]
\centering
\small
\begin{tabular}{@{}ll@{}}
\toprule
\textbf{Component} & \textbf{What it is} \\
\midrule
Machine-readable schema & JSON Schema (draft 2020-12): 43 fields, 19 controlled \\
                        & vocabularies, 4 cross-field conditional rules \\
Pilot dataset           & 47 draft dialogue items \\
Guidelines              & 5 documents: annotation, contrastive construction, \\
                        & adjudication, human evaluation, and training \\
Templates               & Annotation, item-review, and adjudication sheets \\
Tooling                 & Schema validation, dataset validation, statistics, \\
                        & annotation-sheet generation, agreement analysis \\
Tests                   & 41 automated tests \\
\bottomrule
\end{tabular}
\caption{Components produced so far.}
\label{tab:built}
\end{table}

\subsection{The schema, as a contract}
I wrote the schema as a JSON Schema so that validation is enforced by a machine
rather than by discipline. It checks the obvious things, such as required fields
and controlled vocabularies, but also several cross-field rules that catch subtler
mistakes: an item whose identifier contradicts its own labels is rejected, a
code-switching item that fails to name a switch function is rejected, and an item
that references a context turn which does not exist is rejected. Where the design
refined the earlier blueprint, for instance by dropping a politeness level that I
could not reliably separate from its neighbour, I recorded the change and my reason
in a versioned changelog rather than letting the two drift apart silently.

\subsection{The dataset}
The pilot contains 47 items. I aimed for a balanced spread and hit it closely.
Table~\ref{tab:cat} gives the category and language distribution, and
Table~\ref{tab:ctx} the source-type, context, and severity distributions.

\begin{table}[h]
\centering
\small
\begin{minipage}[t]{0.48\textwidth}
\centering
\begin{tabular}{@{}lrr@{}}
\toprule
\textbf{Category} & \textbf{Items} & \textbf{Share} \\
\midrule
Politeness and formality      & 14 & 29.8\% \\
Indirect requests / refusals  & 14 & 29.8\% \\
Stance and emotion            & 12 & 25.5\% \\
Code-switching                & 7  & 14.9\% \\
\midrule
Hindi                         & 26 & 55.3\% \\
Hinglish                      & 21 & 44.7\% \\
\bottomrule
\end{tabular}
\end{minipage}
\hfill
\begin{minipage}[t]{0.48\textwidth}
\centering
\begin{tabular}{@{}lrr@{}}
\toprule
\textbf{Attribute} & \textbf{Items} & \textbf{Share} \\
\midrule
Elicited            & 35 & 74.5\% \\
Minimal pair        & 12 & 25.5\% \\
\midrule
Context required    & 22 & 46.8\% \\
Context helpful     & 17 & 36.2\% \\
Context not required& 8  & 17.0\% \\
\midrule
Severity major      & 37 & 78.7\% \\
Severity minor      & 7  & 14.9\% \\
Severity critical   & 3  & 6.4\% \\
\bottomrule
\end{tabular}
\end{minipage}
\caption{Left: category and language. Right: source type, context necessity, and severity.}
\label{tab:cat}\label{tab:ctx}
\end{table}

Two decisions in this table are deliberate and worth explaining. There are no items
drawn from real corpora, because I could not yet document their provenance or verify
their licensing, and inventing a source would have been fabrication. I
would rather report zero natural items honestly than pad the count. And the twelve
minimal pairs are six matched pairs that differ in exactly one feature, such as the
\code{aap} versus \code{tum} pronoun choice, which lets an evaluation isolate a
single pragmatic contrast.

\subsection{Tooling and tests}
The scripts are the part I am most confident in, because they are testable and
tested. Validation runs in two layers, a JSON Schema pass and a set of cross-field
consistency checks. The statistics script produces the distribution tables above
directly from the data, so the report cannot drift from the dataset. The agreement
script computes percentage agreement, Cohen's kappa, and Krippendorff's alpha, and,
importantly, it reports insufficient data rather than inventing a number when there
is nothing to compute, which is its correct behaviour right now. Forty-one tests
cover the validators and the agreement statistics, including a check of
Krippendorff's alpha against a hand-computed worked example.

% ======================================================================
\section{Reproducibility}
% ======================================================================

Everything can be re-run from the repository root, and I verified all of it. The
commands and their current results are in Table~\ref{tab:verify}.

\begin{table}[h]
\centering
\small
\begin{tabular}{@{}ll@{}}
\toprule
\textbf{Command} & \textbf{Result} \\
\midrule
\code{validate\_schema.py}                  & Schema is a valid draft 2020-12 schema \\
\code{validate\_jsonl.py \dots pilot}        & 47 of 47 items valid \\
\code{python -m unittest discover}          & 41 tests pass \\
\code{generate\_dataset\_statistics.py}      & Distribution as reported above \\
\code{create\_annotation\_sheet.py}          & 94 candidate rows for 47 items \\
\code{calculate\_agreement.py}               & Correctly reports insufficient data \\
\bottomrule
\end{tabular}
\caption{Verification run from the repository root.}
\label{tab:verify}
\end{table}

% ======================================================================
\section{Honest limitations}
% ======================================================================

I think it is more useful to state the weaknesses plainly than to present the work
as finished, so here are the ones I already know about.

\begin{itemize}
  \item \textbf{Nothing is human-validated.} Every item is marked as needing native
  review. No annotation has happened, no agreement has been measured, and I have
  not run any translation systems. The apparatus exists; the measurements do not.
  \item \textbf{Code-switching is kept, and its testability is now something the
  pilot measures.} A Hindi-English switch cannot survive into a monolingual English
  translation, so these items test whether the \emph{effect} of the switch, such as
  a shift into an authoritative or intimate footing, survives. My supervisor and I
  decided to keep the category rather than drop it, and to convert that testability
  worry into a measurement: if annotators cannot agree on code-switch preservation
  while they agree on the other dimensions, that is the evidence the category is not
  well-posed against an English target, and it gets revisited.
  \item \textbf{The target variety is Indian English, and the references do not all
  conform yet.} We decided the benchmark targets Indian English, which matters
  because it judges social meaning in the output. The 47 references were drafted
  before that decision in a general informal register, so bringing them into Indian
  English is now an explicit criterion of the native review rather than something
  already done.
  \item \textbf{Every item has a single author.} I wrote all 47, including both the
  reference and the contrastive translation for each, so any blind spot of mine is
  invisible within the dataset. Independent review is the first thing I want to fix
  next.
  \item \textbf{The severity distribution is skewed} toward \code{MAJOR}, which I
  chose not to correct by relabelling, because rating to a target is exactly the
  calibration error my own guidelines warn against.
\end{itemize}

% ======================================================================
\section{Next steps}
% ======================================================================

The critical path from here does not run through more writing by me; it runs
through people. The immediate steps, in order, are:

\begin{enumerate}
  \item Act on the two decisions now settled with my supervisor: keep the
  code-switching category, and target Indian English. In practice this folds into
  the review pass below, which conforms the references to Indian English and treats
  code-switch agreement as the test of that category.
  \item Recruit two native-speaker annotators and one adjudicator, ideally from
  different regions so that region-dependent judgements can be detected rather than
  hidden.
  \item Run a two-stage pilot: first an item review that removes or fixes any item
  whose contrastive is not genuinely wrong, then blind rating of the survivors.
  \item Measure inter-annotator agreement per dimension and check the central
  assumption, that annotators can hold semantic adequacy and pragmatic preservation
  apart.
  \item Revise the schema in light of the findings, and only then run the first
  translation baselines, since baselines against an unvalidated benchmark would not
  mean anything.
\end{enumerate}

% ======================================================================
\section{Conclusion}
% ======================================================================

I set out to make a specific kind of translation error measurable: the case where
the words are right but the social meaning is wrong. So far I have turned that idea
into a concrete, validated, and reproducible instrument, a schema, a 47-item pilot
dataset, a full set of annotation guidelines, and a tested toolchain, while being
careful not to claim any result that would require human judgement I have not yet
collected. The honest status is that the measuring apparatus is built and the
measurements come next. That is the right place to be before spending scarce
annotator time, and it is what I have set the work up to do.

% ----------------------------------------------------------------------
\begin{thebibliography}{9}
\small

\bibitem{sennrich2016}
R.~Sennrich, B.~Haddow, and A.~Birch.
\newblock Controlling Politeness in Neural Machine Translation via Side Constraints.
\newblock In \emph{NAACL-HLT}, pages 35--40, 2016.

\bibitem{niu2017}
X.~Niu, M.~Martindale, and M.~Carpuat.
\newblock A Study of Style in Machine Translation: Controlling the Formality of
Machine Translation Output.
\newblock In \emph{EMNLP}, pages 2814--2819, 2017.

\bibitem{niu2020}
X.~Niu and M.~Carpuat.
\newblock Controlling Neural Machine Translation Formality with Synthetic Supervision.
\newblock In \emph{AAAI}, pages 8568--8575, 2020.

\bibitem{bawden2018}
R.~Bawden, R.~Sennrich, A.~Birch, and B.~Haddow.
\newblock Evaluating Discourse Phenomena in Neural Machine Translation.
\newblock In \emph{NAACL-HLT}, pages 1304--1313, 2018.

\bibitem{voita2019a}
E.~Voita, R.~Sennrich, and I.~Titov.
\newblock When a Good Translation is Wrong in Context: Context-Aware Machine
Translation Improves on Deixis, Ellipsis, and Lexical Cohesion.
\newblock In \emph{ACL}, pages 1198--1212, 2019.

\bibitem{voita2019b}
E.~Voita, R.~Sennrich, and I.~Titov.
\newblock Context-Aware Monolingual Repair for Neural Machine Translation.
\newblock In \emph{EMNLP-IJCNLP}, pages 877--886, 2019.

\bibitem{agrawal2024}
S.~Agrawal, A.~Farajian, P.~Fernandes, R.~Rei, and A.~F.~T.~Martins.
\newblock Assessing the Role of Context in Chat Translation Evaluation: Is Context
Helpful and Under What Conditions?
\newblock \emph{Transactions of the ACL}, 12:1250--1267, 2024.

\end{thebibliography}

\end{document}
